\documentclass[11pt,a4paper]{article}
\input{../../sablony/pripravy/priprava.sty}
\usepackage{tikz,pgfplots}
\pgfplotsset{compat=1.18}
\nastavPripravu{1. ročník}{Kvadratické rovnice a nerovnice}{103}{Neúplná kvadratická rovnice I}
\begin{document}
\blok{Neúplná kvadratická rovnice I}
Od průsečíku paraboly s osou x k rovnici druhého stupně. Přímé odmocňování a počet reálných řešení.


\begin{definiceBox}
\textbf{Kvadratická rovnice} je rovnice převoditelná ekvivalentními úpravami na
\[
ax^2+bx+c=0,\qquad a,b,c\in\mathbb R,\quad a\ne0.
\]
Koeficient $a$ je \textbf{kvadratický}, $b$ \textbf{lineární} a $c$ \textbf{absolutní člen}.
Kořeny jsou všechna reálná $x$ splňující rovnici; jejich množinu značíme $K$.
\textbf{Úplná kvadratická rovnice} má navíc $b\ne0$ a $c\ne0$.
Tvar $ax^2+bx+c=0$ proto obecně označuje kvadratickou rovnici, nikoli nutně úplnou.
\end{definiceBox}
\begin{definiceBox}
\textbf{Neúplná kvadratická rovnice} má alespoň jeden z koeficientů $b,c$ roven nule.
Rozlišujeme tvary $ax^2+bx=0$ a $ax^2+c=0$ (včetně případu $ax^2=0$).
\end{definiceBox}
\poznamka{Terminologie \uv{ryze kvadratická} není jednotná. Často se tím myslí $ax^2+c=0$ ($b=0$); případ $ax^2=0$ má navíc $c=0$. Pro žáky je důležité rozpoznat tvar a správnou metodu, nikoli memorovat sporné názvosloví.}
\blok{Po odmocnění musí zůstat absolutní hodnota}
Pro reálná čísla platí $\sqrt{x^2}=|x|$. Odtud plyne
\[
\begin{aligned}
x^2&=9\
|x|&=3\
x&=-3\quad\text{nebo}\quad x=3,\qquad K=\{-3;3\}.
\end{aligned}
\]
\textbf{Pozor:} $x^2=-9$ nemá reálné řešení; $x^2=0$ má jediný kořen. Stejnou zásadu
použijeme i při odmocňování výrazu $(x-p)^2$.
\blok{Tvar nemusí být hned uspořádán podle mocnin}
Na výsledný \emph{nulový tvar} přivedeme všechny členy k jedné straně a zapíšeme
je sestupně: nejprve $x^2$, poté $x$, nakonec absolutní člen.
Například $5=2x^2+x$ upravíme na $2x^2+x-5=0$.

\blok{Společně řešené příklady}
\begin{prikladReseniBox}[title={Příklad 1}]
$x^{2} - 16=0$
\tcblower
\[
\begin{aligned}
x^2&=16\\|x|&=4\\[2pt]
K&=\{-4;4\}
\end{aligned}
\]
\end{prikladReseniBox}
\begin{prikladReseniBox}[title={Příklad 2}]
$2x^{2} + 8=0$
\tcblower
\[
\begin{aligned}
x^2&=-4\\[2pt]
K&=\varnothing
\end{aligned}
\]
\end{prikladReseniBox}
\begin{prikladReseniBox}[title={Příklad 3}]
$3x^{2}=0$
\tcblower
\[
\begin{aligned}
3x^2&=0\ \Rightarrow\ x=0\\[2pt]
K&=\{0\}
\end{aligned}
\]
\end{prikladReseniBox}
\begin{prikladReseniBox}[title={Příklad 4 – nejprve uprav na nulový tvar}]
$4x^2=12x$.
\tcblower
\[\begin{aligned}
4x^2-12x&=0\\4x(x-3)&=0\\K&=\{0;3\}
\end{aligned}\]
\end{prikladReseniBox}
\blok{Samostatné procvičování}
\begin{zadaniReseniBox}[title={Úloha 1}]
$x^{2} - 9=0$
\tcblower
\[
\begin{aligned}
x^2&=9\\|x|&=3\\[2pt]
K&=\{-3;3\}
\end{aligned}
\]
\end{zadaniReseniBox}
\begin{zadaniReseniBox}[title={Úloha 2}]
$2x^{2} - 18=0$
\tcblower
\[
\begin{aligned}
x^2&=9\\|x|&=3\\[2pt]
K&=\{-3;3\}
\end{aligned}
\]
\end{zadaniReseniBox}
\blok{Volitelná rezerva}
\poznamka{Studenti si vymyslí obdobný příklad se změnou jednoho koeficientu a vysvětlí, jak se změní počet kořenů či řešení.}
\blok{Rozšiřující příklady ze sbírky}
\poznamka{Ukázka ryze neúplné rovnice s necelými kořeny: stihneme-li, porovnej s grafem.}

\begin{zadaniReseniBox}[title={Rozšiřující příklad}]
Řešte v $\mathbb R$: $2x^{2} - 5=0$.
\tcblower
\[
\begin{aligned}2x^2-5&=0\\x^2&=\frac52\\|x|&=\sqrt{\frac52}\\K&=\{-\frac{\sqrt{10}}2;\frac{\sqrt{10}}2\}\end{aligned}
\]
\end{zadaniReseniBox}

\blok{Domácí úkol}
\begin{zadaniReseniBox}[title={Domácí příprava}]
\begin{enumerate}[label=\alph*),itemsep=2pt]
\item $x^{2} - 25=0$
\item $3x^{2} - 27=0$
\item $2x^{2} + 8=0$
\item $5x^{2}=0$
\item $2x^{2} - 7=0$
\item $- 3x^{2} - 12=0$
\end{enumerate}
\tcblower
\begin{enumerate}[label=\alph*),itemsep=2pt]
\item
\[\begin{aligned}
x^2&=25\\|x|&=5\\[2pt]
K&=\{-5;5\}
\end{aligned}\]
\item
\[\begin{aligned}
x^2&=9\\|x|&=3\\[2pt]
K&=\{-3;3\}
\end{aligned}\]
\item
\[\begin{aligned}
x^2&=-4\\[2pt]
K&=\varnothing
\end{aligned}\]
\item
\[\begin{aligned}
3x^2&=0\ \Rightarrow\ x=0\\[2pt]
K&=\{0\}
\end{aligned}\]
\item
\[\begin{aligned}
x^2&=\frac{7}{2}\\|x|&=\sqrt{\frac72}\\[2pt]
K&=\{- \frac{\sqrt{14}}{2};\frac{\sqrt{14}}{2}\}
\end{aligned}\]
\item
\[\begin{aligned}
x^2&=-4\\[2pt]
K&=\varnothing
\end{aligned}\]
\end{enumerate}
\end{zadaniReseniBox}
\poznamka{Orientační čas: úvod 5 min, společné příklady 18 min, samostatně 15 min, kontrola a shrnutí 7 min. Rezerva je dobrovolná. Číselné parametry modelových aplikací nejsou skutečně naměřená data.}
\end{document}
